\subsection{The Exponential and Gamma Distributions}
\hrulefill

\subsubsection*{The Exponential Distribution}
A continuous random variable $X$ is said to have an Exponential distribution with parameter
$\lambda > 0$ if its probability density function is given by
\[f(x) = \begin{cases}
\lambda e^{-\lambda x} & x > 0 \\
0 & \text{otherwise}
\end{cases}.\]

\paragraph*{Notation:} $X \sim \text{Exp}(\lambda)$ means that $X$ is an Exponential random variable with parameter
$\lambda$.

\paragraph*{CDF:} The CDF of $X \sim \text{Exp}(\lambda)$ is
\[F(x) = P(X \leq x) = \begin{cases}
1 - e^{-\lambda x} & x \geq 0 \\
0 & x < 0
\end{cases}.\]

\paragraph*{Used for:}
\begin{itemize}
    \item The time until a radioactive particle decays.
    \item The time until a customer arrives at a service point.
    \item The time until a light bulb burns out.
    \item Life spans of living organisms.
\end{itemize}

\paragraph*{Survival Function:} The survival function is given by
\[S(t) = P(X > t) = 1 - F(t) = e^{-\lambda t}.\]
This is the probability that the event of interest has not occurred by time $t$.
This is done instead of the CDF for lifespans because people are little bitches about death, so now we have to 
learn more formulae.

\paragraph*{Mean and Variance:} If $X \sim \text{Exp}(\lambda)$, then
\[\mu = E(X) = \frac{1}{\lambda}, \quad \sigma^2 = Var(X) = \frac{1}{\lambda^2}.\]

<<<<<<< HEAD
\paragraph*{Note:} When integrating by parts, let $u$ be the polynomial term and $dv$ be the exponential term.
=======
\paragraph*{Note:} When integrating by parts, let $u$ be the polynomial term and $dv$ be the exponential term.
\paragraph*{Note:} $\lim_{b\to\infty} \frac{\text{Polynomial}}{e^{\lambda b}} = 0$ for any polynomial.
>>>>>>> 5595871 (Stats - Class Notes 10/15)
